\chapter[Band 39: Halbringe]{Band 39 auf einen Blick}
\noindent\textbf{Halbringe}\par\medskip
Die additiven und multiplikativen Monoidaxiome werden hier verbunden.
Ein Halbring muss nicht \(0\ne1\) erfüllen; die Multiplikation muss
nicht kommutativ sein. Die natürliche Abbildung wird mittels der bereits
entwickelten Iteration konstruiert.

\begin{BandUebersicht}
\textbf{Halbringaxiome}
Ein Halbring \((R,+,\cdot,0,1)\) erfüllt:\UebFormel{\begin{aligned}
&\CommutativeMonoidAlg R+0,\\
&\MonoidAlg R\cdot1,\\
&x(y+z)=xy+xz,\\
&(x+y)z=xz+yz,\\
&0x=0=x0\qquad(x,y,z\in R).
\end{aligned}}Kommutativ heißt zusätzlich \(xy=yx\) für \(x,y\in R\).
\UebQuellen{\FormulaRefAuto{SemiringDef}[def];
\FormulaRefAuto{CommutativeUnitalSemiringDef}[def]}
&
\textbf{Natürliche Zahlen als Halbring}
\UebFormel{\begin{aligned}
&\SemiringAlg{\mathbb N}{+}{\cdot}{0}{1},\\
&\CommutativeSemiringAlg{\mathbb N}{+}{\cdot}{0}{1}.
\end{aligned}}Die additive Inversenbedingung eines Rings gilt hier nicht:
insbesondere hat \(1\) kein additives Inverses in \(\mathbb N\).
\UebQuellen{\FormulaRefAuto{NaturalUnitalSemiring}[thm];
\FormulaRefAuto{NaturalCommutativeUnitalSemiring}[thm];
\FormulaRefAuto{NaturalNumbersNotRing}[thm]} \\
\addlinespace[.8ex]

\textbf{Morphismen}
Für Halbringe \(\mathfrak R,\mathfrak S\) ist ein
Halbringhomomorphismus \(f:R\to S\) durch\UebFormel{\begin{aligned}
&f(x+y)=f(x)+f(y),\\
&f(xy)=f(x)f(y),\\
&f(0_R)=0_S,\qquad f(1_R)=1_S
\end{aligned}}gegeben. Ein Isomorphismus ist ein bijektiver solcher Morphismus.
\UebQuellen{\FormulaRefAuto{SemiringHomDef}[def];
\FormulaRefAuto{SemiringIsoDef}[def]}
&
\textbf{Komposition}
Für Halbringe \(\mathfrak R,\mathfrak S,\mathfrak T\):\UebFormel{\begin{aligned}
&\SemiringHom f{\mathfrak R}{\mathfrak S},\\
&\SemiringHom g{\mathfrak S}{\mathfrak T}\\
&\qquad\vdash
 \SemiringHom{g\circ f}{\mathfrak R}{\mathfrak T}.
\end{aligned}}
\UebQuellen{\FormulaRefAuto{SemiringHomComposition}[thm]} \\
\addlinespace[.8ex]

\textbf{Die natürliche Abbildung}
Für einen Halbring
\(\mathfrak R=(R,\oplus,\odot,0_R,1_R)\):\UebFormel{\begin{aligned}
&\sigma_R(x)=x\oplus1_R,\\
&\eta_R(0)=0_R,\\
&\eta_R(n+1)=\sigma_R(\eta_R(n))
 \quad(n\in\mathbb N).
\end{aligned}}
\UebQuellen{\FormulaRefAuto{SemiringSuccessorDef}[def];
\FormulaRefAuto{NaturalSemiringMapDef}[def]}
&
\textbf{Universelle Eigenschaft}
Die Abbildung \(\eta_R:\mathbb N\to R\) erfüllt für \(m,n\in\mathbb N\):\UebFormel{\begin{aligned}
&\eta_R(1)=1_R,\\
&\eta_R(m+n)=\eta_R(m)\oplus\eta_R(n),\\
&\eta_R(mn)=\eta_R(m)\odot\eta_R(n).
\end{aligned}}Sie ist der eindeutig bestimmte Halbringhomomorphismus
von \((\mathbb N,+,\cdot,0,1)\) nach \(\mathfrak R\).
\UebQuellen{\FormulaRefAuto{NaturalSemiringMapOne}[thm];
\FormulaRefAuto{NaturalSemiringMapAddition}[thm];
\FormulaRefAuto{NaturalSemiringMapMultiplication}[thm];
\FormulaRefAuto{NaturalSemiringUniversalHomomorphism}[thm]} \\
\addlinespace[.8ex]

\textbf{Nachfolgerinduktion im Ziel}
Die Induktionsbedingung lautet:\UebFormel{\begin{aligned}
&A\subseteq R,\quad 0_R\in A,\\
&\forall x\in A\ (\sigma_R(x)\in A)\\
&\qquad\Longrightarrow A=R.
\end{aligned}}
\UebQuellen{\FormulaRefAuto{NaturalSemiringMapSurjectiveIffSuccessorInduction}[thm]}
&
\textbf{Surjektivität und Injektivität}
Für einen Halbring \(R\) ist diese Bedingung äquivalent zu
\(\eta_R:\mathbb N\sur R\). Außerdem gilt:\UebFormel{\begin{aligned}
&\sigma_R:R\inj R,\quad
 \forall x\in R\ (\sigma_R(x)\ne0_R)\\
&\qquad\vdash\eta_R:\mathbb N\inj R.
\end{aligned}}
\UebQuellen{\FormulaRefAuto{NaturalSemiringMapSurjectiveIffSuccessorInduction}[thm];
\FormulaRefAuto{NaturalSemiringMapInjectivePeanoCriterion}[thm]} \\
\addlinespace[.8ex]

\textbf{Endlichkeit und Isomorphie}
Ein endlicher Halbring erfüllt zusätzlich
\(\Finite R\). Die kanonische Abbildung bleibt stets \(\eta_R\).
\UebQuellen{\FormulaRefAuto{FiniteSemiringDef}[def];
\FormulaRefAuto{NaturalSemiringMapDef}[def]}
&
\textbf{Grenzen und exaktes Kriterium}
\UebFormel{\begin{aligned}
&\SemiringAlg R\oplus\odot{0_R}{1_R},
 \ \Finite R\\
&\qquad\vdash\neg(\eta_R:\mathbb N\inj R).
\end{aligned}}Für beliebige Halbringe ist \(\eta_R\) genau dann ein
Halbringisomorphismus, wenn es injektiv und surjektiv ist.
\UebQuellen{\FormulaRefAuto{NaturalSemiringMapNotInjectiveIntoFiniteSemiring}[thm];
\FormulaRefAuto{NaturalSemiringMapIsoIffInjectiveAndSurjective}[thm]} \\
\addlinespace[.8ex]
\end{BandUebersicht}
